% GENERATED FILE. Edit canonical lessons, then run npm run generate:books.
% canonical-source-hash: fnv1a64:48485d01b8f3b58a
% canonical-lessons: TE-C547-subhraparacu, TE-C547-utuku, TE-C547-araveyu, TE-C547-dulupu, TE-C547-tagilincu

\chapter{Housework}
\label{ch:te-housework}
\hlchaptermodality{547}

\begin{chapteropening}
\textbf{By the end of this chapter:} \emph{I can talk about cleaning, washing, drying, dusting and hanging things up.}

Complete the last lesson of chapter 547: I can talk about cleaning, washing, drying, dusting and hanging things up.
\end{chapteropening}

\section[\texorpdfstring{\te{శుభ్రపరచు} (śubhraparacu)}{శుభ్రపరచు (śubhraparacu)}]{\te{శుభ్రపరచు} (śubhraparacu) — to cleanse; to clean}

\label{lesson:TE-C547-subhraparacu}



\begin{quote}
Before the new one: say the Telugu for a small water pot, then the Telugu for a cockroach.
\end{quote}

\subsection*{You'll want to know: \te{శుభ్రపరచు}}
\textbf{\te{శుభ్రపరచు}} — \emph{śubhraparacu} — \textquotedblleft{}to cleanse; to clean\textquotedblright{}.

\begin{quote}
\te{పండుగకు} \te{ముందు} \te{ఇల్లంతా} \te{శుభ్రపరిచాము}.
\end{quote}

\begin{quote}
\emph{paṇḍugaku mundu illantā śubhraparicāmu.}
\end{quote}

\begin{quote}
Before the festival we cleaned the whole house.
\end{quote}

Past (I …): \te{శుభ్రపరిచాను}, \emph{śubhraparicānu}.

Formal; \te{శుభ్రం} \te{చేయు} in speech.

\begin{cousinweb}
\textbf{\te{శుభ్రపరచు}} (\emph{śubhraparacu}) is built from \te{శుభ్రం}, \textquotedblleft{}clean\textquotedblright{}, from Sanskrit \emph{śubhra}, \textquotedblleft{}bright, white\textquotedblright{} (CDIAL 12539) and \emph{paracu}, \textquotedblleft{}to cause, to make\textquotedblright{}, as in \te{ఆశ్చర్యపరచు}.
\end{cousinweb}

\begin{grammarlens}[title={What you've built}]
One more word for describing what you lived and felt.
\end{grammarlens}

\subsection*{Your turn}
\begin{itemize}
\raggedright
  \item \emph{Say it:} \emph{śubhraparacu}
  \item \emph{Say it:} \emph{śubhraparacu}, once more
  \item \emph{Say it:} \emph{paṇḍugaku mundu illantā śubhraparicāmu.} — Before the festival we cleaned the whole house.
  \item \emph{Recall:} say the Telugu for a small water pot, then the Telugu for a cockroach, then say \emph{śubhraparacu} again
\end{itemize}

\begin{tcolorbox}[breakable,colback=teal!4,colframe=teal!35!black,title={Before you move on}]
What does \te{శుభ్రపరచు} mean? (\textquotedblleft{}To cleanse; to clean\textquotedblright{}.) Say it once more.
\end{tcolorbox}
\section[\texorpdfstring{\te{ఉతుకు} (utuku)}{ఉతుకు (utuku)}]{\te{ఉతుకు} (utuku) — to wash clothes}

\label{lesson:TE-C547-utuku}



\begin{quote}
Before the new one: say the Telugu for a cockroach, then the Telugu for to cleanse.
\end{quote}

\subsection*{You'll want to know: \te{ఉతుకు}}
\textbf{\te{ఉతుకు}} — \emph{utuku} — \textquotedblleft{}to wash clothes\textquotedblright{}.

\begin{quote}
\te{ఆదివారం} \te{నేను} \te{బట్టలు} \te{ఉతుకుతాను}.
\end{quote}

\begin{quote}
\emph{ādivāraṁ nēnu baṭṭalu utukutānu.}
\end{quote}

\begin{quote}
On Sundays I wash the clothes.
\end{quote}

Past (I …): \te{ఉతికాను}, \emph{utikānu}.

For clothes; \te{కడుగు} is for dishes, hands and floors.

\begin{cousinweb}
\textbf{\te{ఉతుకు}} (\emph{utuku}) is native Dravidian: DEDR 2667 lists Telugu \emph{utuku}, \textquotedblleft{}to wash, as clothes\textquotedblright{}, with Kolami \emph{sutk-}.
\end{cousinweb}

\begin{grammarlens}[title={What you've built}]
One more word for describing what you lived and felt.
\end{grammarlens}

\subsection*{Your turn}
\begin{itemize}
\raggedright
  \item \emph{Say it:} \emph{utuku}
  \item \emph{Say it:} \emph{utuku}, once more
  \item \emph{Say it:} \emph{ādivāraṁ nēnu baṭṭalu utukutānu.} — On Sundays I wash the clothes.
  \item \emph{Recall:} say the Telugu for a cockroach, then the Telugu for to cleanse, then say \emph{utuku} again
\end{itemize}

\begin{tcolorbox}[breakable,colback=teal!4,colframe=teal!35!black,title={Before you move on}]
What does \te{ఉతుకు} mean? (\textquotedblleft{}To wash clothes\textquotedblright{}.) Say it once more.
\end{tcolorbox}
\section[\texorpdfstring{\te{ఆరవేయు} (āravēyu)}{ఆరవేయు (āravēyu)}]{\te{ఆరవేయు} (āravēyu) — to hang to dry; to put out to dry}

\label{lesson:TE-C547-araveyu}



\begin{quote}
Before the new one: say the Telugu for to cleanse, then the Telugu for to wash clothes.
\end{quote}

\subsection*{You'll want to know: \te{ఆరవేయు}}
\textbf{\te{ఆరవేయు}} — \emph{āravēyu} — \textquotedblleft{}to hang to dry; to put out to dry\textquotedblright{}.

\begin{quote}
\te{ఉతికిన} \te{బట్టలను} \te{ఎండలో} \te{ఆరవేశాను}.
\end{quote}

\begin{quote}
\emph{utikina baṭṭalanu eṇḍalō āravēśānu.}
\end{quote}

\begin{quote}
I hung the washed clothes out in the sun to dry.
\end{quote}

Past (I …): \te{ఆరవేశాను}, \emph{āravēśānu}.

\te{ఆరు}, to dry, + \te{వేయు}, to put.

\begin{cousinweb}
\textbf{\te{ఆరవేయు}} (\emph{āravēyu}) is native Dravidian: DEDR 404 lists Telugu \emph{āṟu}, \textquotedblleft{}to cool; to dry, to become dry\textquotedblright{}; \emph{vēyu}, \textquotedblleft{}to put\textquotedblright{}, makes it something you do to the clothes.
\end{cousinweb}

\begin{grammarlens}[title={What you've built}]
One more word for describing what you lived and felt.
\end{grammarlens}

\subsection*{Your turn}
\begin{itemize}
\raggedright
  \item \emph{Say it:} \emph{āravēyu}
  \item \emph{Say it:} \emph{āravēyu}, once more
  \item \emph{Say it:} \emph{utikina baṭṭalanu eṇḍalō āravēśānu.} — I hung the washed clothes out in the sun to dry.
  \item \emph{Recall:} say the Telugu for to cleanse, then the Telugu for to wash clothes, then say \emph{āravēyu} again
\end{itemize}

\begin{tcolorbox}[breakable,colback=teal!4,colframe=teal!35!black,title={Before you move on}]
What does \te{ఆరవేయు} mean? (\textquotedblleft{}To hang to dry; to put out to dry\textquotedblright{}.) Say it once more.
\end{tcolorbox}
\section[\texorpdfstring{\te{దులుపు} (dulupu)}{దులుపు (dulupu)}]{\te{దులుపు} (dulupu) — to shake out; to dust}

\label{lesson:TE-C547-dulupu}



\begin{quote}
Before the new one: say the Telugu for to wash clothes, then the Telugu for to hang to dry.
\end{quote}

\subsection*{You'll want to know: \te{దులుపు}}
\textbf{\te{దులుపు}} — \emph{dulupu} — \textquotedblleft{}to shake out; to dust\textquotedblright{}.

\begin{quote}
\te{అమ్మ} \te{దుప్పట్లు} \te{దులిపి} \te{ఎండలో} \te{వేసింది}.
\end{quote}

\begin{quote}
\emph{amma duppaṭlu dulipi eṇḍalō vēsindi.}
\end{quote}

\begin{quote}
Mum shook out the blankets and put them in the sun.
\end{quote}

Past (I …): \te{దులిపాను}, \emph{dulipānu}.

Shaking the dust off cloth or a mat.

\begin{cousinweb}
\textbf{\te{దులుపు}} (\emph{dulupu}) is native Dravidian: DEDR 3359 lists Telugu \emph{dulupu}, \textquotedblleft{}to shake so as to remove dust, to shake off\textquotedblright{}, in the article of Tamil \emph{tuḷaṅku}, \textquotedblleft{}to move, to shake\textquotedblright{}.
\end{cousinweb}

\begin{grammarlens}[title={What you've built}]
One more word for describing what you lived and felt.
\end{grammarlens}

\subsection*{Your turn}
\begin{itemize}
\raggedright
  \item \emph{Say it:} \emph{dulupu}
  \item \emph{Say it:} \emph{dulupu}, once more
  \item \emph{Say it:} \emph{amma duppaṭlu dulipi eṇḍalō vēsindi.} — Mum shook out the blankets and put them in the sun.
  \item \emph{Recall:} say the Telugu for to wash clothes, then the Telugu for to hang to dry, then say \emph{dulupu} again
\end{itemize}

\begin{tcolorbox}[breakable,colback=teal!4,colframe=teal!35!black,title={Before you move on}]
What does \te{దులుపు} mean? (\textquotedblleft{}To shake out; to dust\textquotedblright{}.) Say it once more.
\end{tcolorbox}
\section[\texorpdfstring{\te{తగిలించు} (tagiliñcu)}{తగిలించు (tagiliñcu)}]{\te{తగిలించు} (tagiliñcu) — to hang up; to hook on}

\label{lesson:TE-C547-tagilincu}



\begin{quote}
Before the new one: say the Telugu for to hang to dry, then the Telugu for to shake out.
\end{quote}

\subsection*{You'll want to know: \te{తగిలించు}}
\textbf{\te{తగిలించు}} — \emph{tagiliñcu} — \textquotedblleft{}to hang up; to hook on\textquotedblright{}.

\begin{quote}
\te{అతను} \te{చొక్కాను} \te{మేకుకు} \te{తగిలించాడు}.
\end{quote}

\begin{quote}
\emph{atanu cokkānu mēkuku tagiliñcāḍu.}
\end{quote}

\begin{quote}
He hung his shirt on a nail.
\end{quote}

Past (I …): \te{తగిలించాను}, \emph{tagiliñcānu}.

-\te{కి} \te{తగిలించు}, to hang on a peg or a nail.

\begin{cousinweb}
\textbf{\te{తగిలించు}} (\emph{tagiliñcu}) is native Dravidian: DEDR 3004 lists Telugu \emph{tagulu}, \textquotedblleft{}to touch, to catch\textquotedblright{}, and \emph{tagilincu}, \textquotedblleft{}to cause to touch, to catch or to stick\textquotedblright{}.
\end{cousinweb}

\begin{grammarlens}[title={What you've built}]
One more word for describing what you lived and felt.
\end{grammarlens}

\subsection*{Your turn}
\begin{itemize}
\raggedright
  \item \emph{Say it:} \emph{tagiliñcu}
  \item \emph{Say it:} \emph{tagiliñcu}, once more
  \item \emph{Say it:} \emph{atanu cokkānu mēkuku tagiliñcāḍu.} — He hung his shirt on a nail.
  \item \emph{Recall:} say the Telugu for to hang to dry, then the Telugu for to shake out, then say \emph{tagiliñcu} again
\end{itemize}

\begin{tcolorbox}[breakable,colback=teal!4,colframe=teal!35!black,title={Before you move on}]
What does \te{తగిలించు} mean? (\textquotedblleft{}To hang up; to hook on\textquotedblright{}.) Say it once more.
\end{tcolorbox}
